\section{El manuscrito como demostrador reproducible}
El repositorio que acompaña a este trabajo aplica el enfoque a la producción de los propios artículos. Cada manuscrito dispone de carpetas para texto, datos, fuentes, figuras y PDF. LaTeX reúne el contenido y las referencias; las figuras se generan como archivos vectoriales independientes. Un mismo comando realiza las comprobaciones y la compilación en el entorno de edición y en integración continua.

La configuración de GitHub Actions separa compilación y escritura. Las propuestas de cambio se compilan con permisos de lectura. En la rama predeterminada, un segundo trabajo recupera las salidas y actualiza los PDF mediante un commit. El token incorporado de GitHub evita que ese push dispare recursivamente otro flujo de push \cite{actions}. Una comprobación de revisión impide publicar un PDF generado a partir de fuentes anteriores a la cabeza actual de la rama. No se fuerza la escritura si la rama cambió.

Los nombres de los PDF son estables y cada nueva compilación sustituye los archivos visibles. Git conserva el historial, de modo que una única ruta actual no implica borrar versiones anteriores. Un manifiesto relaciona el PDF con hashes de sus fuentes y registra el motor de compilación. Este mecanismo facilita rastrear el producto editorial sin prometer identidad binaria entre versiones diferentes de TeX o paquetes.

El demostrador amplía un ejercicio previo para el Instituto de Ingeniería de Venezuela, cuyo repositorio separa datos, análisis, figuras e informes \cite{instituto}. Allí las automatizaciones entregan artifacts; aquí se añade la actualización controlada de los manuscritos en el repositorio. La diferencia responde al objetivo editorial de que un lector encuentre siempre el PDF actual. No se presenta la infraestructura de este paquete como un servicio ya desplegado en GitHub: requiere importar el repositorio y realizar su primera ejecución remota.

La verificación automática comprueba límites de los resúmenes, orden de referencias y errores de compilación. La revisión visual y científica continúa siendo necesaria. Un PDF actualizado puede contener una afirmación equivocada; la trazabilidad ayuda a localizarla y corregirla, pero no reemplaza la responsabilidad de los autores. La apertura pública del borrador debe acordarse con la política de la revista destinataria.
